% TOP_PROBLEM: {"id": 1283, "title": "Transcendence of Apéry's Constant", "category_id": 1, "category": {"id": 1, "name": "number_theory", "display_name": "Number Theory", "description": "Properties of integers, prime numbers, Diophantine equations.", "slug": "number-theory", "order_index": 1, "created_at": "2026-07-31T15:26:25.670Z"}, "status": "open"}
% FIRST_POSTED: 2026-09-27
% !TeX program = lualatex
% Compile twice: lualatex open_problems_500.tex
% All references and prose are embedded; no BibTeX database or external inputs.
\documentclass[11pt,a4paper]{article}
\usepackage[margin=24mm,headheight=14pt]{geometry}
\usepackage{fontspec}
\setmainfont{Latin Modern Roman}
\setsansfont{Latin Modern Sans}
\setmonofont{Latin Modern Mono}
\usepackage{amsmath,amssymb,mathtools}
\usepackage{unicode-math}
\setmathfont{Latin Modern Math}
\usepackage{microtype}
\usepackage{enumitem,xurl,needspace}
\usepackage[dvipsnames]{xcolor}
\usepackage{hyperref}
\hypersetup{unicode=true,colorlinks=true,linkcolor=MidnightBlue,urlcolor=MidnightBlue,
 pdftitle={500 Mathematical Problem Statements},pdfauthor={Source-annotated compilation},bookmarksnumbered=true}
\usepackage{bookmark}
\usepackage{tocloft}
\setlength{\cftsecnumwidth}{3em}
\cftsetpnumwidth{2.5em}
\cftsetrmarg{4em}
\usepackage{titlesec}
\titleformat{\section}{\Large\bfseries\raggedright}{\thesection}{0.7em}{}
\usepackage{fancyhdr}
\pagestyle{fancy}\fancyhf{}
\fancyhead[L]{\small\sffamily Mathematical problem statements}
\fancyhead[R]{\small Source edition 22}
\fancyfoot[C]{\thepage}
\setlength{\parindent}{0pt}
\setlength{\parskip}{5pt plus 1pt}
\setlength{\emergencystretch}{3em}
\setcounter{tocdepth}{1}
\setcounter{secnumdepth}{1}
\setlist[itemize]{leftmargin=3em,itemsep=5pt,topsep=4pt,parsep=0pt}
\allowdisplaybreaks
\newcommand{\N}{\mathbb{N}}
\newcommand{\Z}{\mathbb{Z}}
\newcommand{\Q}{\mathbb{Q}}
\newcommand{\R}{\mathbb{R}}
\newcommand{\C}{\mathbb{C}}
\newcommand{\F}{\mathbb{F}}
\newcommand{\A}{\mathbb{A}}
\newcommand{\PP}{\mathbb P}
\newcommand{\E}{\mathbb{E}}
\newcommand{\eps}{\varepsilon}
\newcommand{\dd}{\,\mathrm d}
\newcommand{\sref}[2]{\hyperlink{source:#1:#2}{[S#2]}}
\newcommand{\eref}[2]{\hyperlink{extra:#1:#2}{[E#2]}}
% Turn the next switch off for a shorter reading copy. The quotations preserve
% every exactTarget field, including qualifications omitted from a short summary.
\newif\ifshowcatalogue
\showcataloguetrue
\newcommand{\cataloguescope}[1]{\ifshowcatalogue
 \paragraph{Exact scope in the supplied catalogue.}
 \begin{quote}\small #1\end{quote}\fi}

% Compile twice with: lualatex 1283.tex
\numberwithin{equation}{section}
\hypersetup{pdftitle={Research attempt -- problem 1283},pdfauthor={Research notebook; original compilation retained}}
\fancyhead[L]{\small\sffamily Research notebook}
\fancyhead[R]{\small\sffamily Problem 1283 | 27 September 2026}
\begin{document}
\setcounter{section}{0}
% ================================================================
% problemId: 1283
% Canonical problem ID: 1283
% exactTarget SHA-256: 5a6860e14a77044631564aeef02ac54fc78e4da15bef48a3ab24544e48cb929a
\clearpage
\section*{\texorpdfstring{Transcendence of Apéry's Constant}{Transcendence of Apéry's Constant}}\label{1283}
\subsection{Definitions and mathematical statement}
Ap\'ery's constant is the absolutely convergent series
\[
 \zeta(3)=\sum_{n=1}^{\infty}\frac1{n^3}.
\]
A real or complex number $z$ is algebraic over $\Q$ if some nonzero polynomial in $\Q[T]$ vanishes at $z$, and transcendental otherwise. Thus the selected assertion is
$P(\zeta(3))\ne0$ for every $P\in\Q[T]\setminus\{0\}$.
A negative resolution would give algebraicity, not necessarily rationality. Irrationality excludes only polynomial relations of degree one; it does not by itself settle this transcendence question. The target concerns this single value, not algebraic independence of a collection of odd zeta values.
\par\smallskip\textit{Formulation and concept references:} \sref{1283}{1}.
\cataloguescope{Prove or disprove that the value zeta(3) is transcendental over the rational numbers.}
\subsection{Short English statement}
Is the sum of the reciprocal cubes of the positive integers transcendental, rather than merely irrational?
% BEGIN RESEARCH ADDITION ATTEMPT-184
\subsection{Research attempt}

\paragraph{Current frontier.}
Ap\'ery's irrationality theorem does not exclude an algebraic value of degree two or higher. The linear-form methods in \sref{1283}{1} address arithmetic of odd zeta values, not transcendence of this particular value. Cohen--Zudilin's 2025 analysis revisits the recurrence and its analytic interpolation; it does not assert the required transcendence theorem \eref{1283}{1}. The catalogue's status label is not a proof: no purported transcendence proof is used here. The calculation below takes the classical recurrence, its limit, and its exponential asymptotics as established inputs and tests whether powering its linear forms can supply the missing higher-degree information.

\paragraph{Attack route.}
There are three genuinely different requirements. A rational approximation must beat a denominator bound to exclude rationality; a degree-$d$ algebraic exclusion needs a bound calibrated to $d$; and a transcendence argument must work for every fixed $d$. One possible route is a new family of simultaneous approximants to $1,\zeta(3),\ldots,\zeta(3)^d$, with a determinant nonvanishing proof. A less expensive proposal is to tensor the classical approximants. We investigate the latter first, while extracting a rigorous interval certificate that does not depend on numerical recognition of the constant.

\paragraph{Attempt.}
Set $a_0=1,a_1=5,b_0=0,b_1=6$, and let both sequences obey
\begin{equation}\label{eq184:rec}
(n+1)^3y_{n+1}=(34n^3+51n^2+27n+5)y_n-n^3y_{n-1}\qquad(n\ge1).
\end{equation}
The established Ap\'ery facts used here are
\[
 a_n=\sum_{j=0}^n\binom nj^2\binom{n+j}j^2\in\Z,
 \quad D_n^3b_n\in\Z,\quad b_n/a_n\longrightarrow\zeta(3),
 \qquad D_n=\operatorname{lcm}(1,\ldots,n).
\]
The recurrence, normalization, and convergence are discussed in \eref{1283}{1}, Section~1. We now derive an enclosure directly from them.

\textbf{Lemma (exact enclosure).} For every $n\ge0$,
\begin{equation}\label{eq184:interval}
\frac{b_n}{a_n}<\zeta(3)
 \le \frac{b_n}{a_n}+\frac{25}{4(n+1)^3a_na_{n+1}}.
\end{equation}
Indeed, the Casoratian $W_n=a_nb_{n+1}-a_{n+1}b_n$ satisfies
$W_n=(n/(n+1))^3W_{n-1}$ and $W_0=6$. Consequently
\[
 \frac{b_{n+1}}{a_{n+1}}-\frac{b_n}{a_n}
 =\frac{6}{(n+1)^3a_na_{n+1}}=:t_n>0.
\]
Also $a_{n+1}\ge5a_n$. To verify the induction, use
$a_{n-1}\le a_n/5$ in \eqref{eq184:rec}; the numerator left after subtracting $5(n+1)^3a_n$ is
$((144/5)n^3+36n^2+12n)a_n\ge0$.
Thus $t_{n+1}/t_n\le1/25$. Summing the positive tail, justified by convergence of $b_n/a_n$, gives
$\zeta(3)-b_n/a_n=\sum_{j\ge n}t_j\le(25/24)t_n$, proving the lemma. Notice that no unproved asymptotic or floating-point comparison enters this enclosure.

For a prescribed polynomial $P$, rational interval arithmetic on \eqref{eq184:interval} can sometimes certify $P(\zeta(3))\ne0$. It is not an algorithm known to terminate for \emph{all} nonzero $P$: failure to exclude zero could persist precisely at an algebraic relation. Using $n=12$, the accompanying script certifies
\[
 |P(\zeta(3))|>2^{-24}
\]
for all $19{,}682$ nonzero polynomials of degree at most eight with
coefficients in $\{-1,0,1\}$. It uses a common integer denominator for
all interval endpoints, not floating-point evaluation. It also checks
the Casoratian and denominator integrality through index 20. These
are bounded certificates, not a transcendence criterion.

Next put $\rho=17+12\sqrt2$, $A_n=D_n^3a_n$, $B_n=D_n^3b_n$. The classical asymptotics give
\begin{align*}
 \log\max(|A_n|,|B_n|)&=(3+\log\rho)n+o(n),\\
 \log|A_n\zeta(3)-B_n|&=-(\log\rho-3)n+o(n).
\end{align*}
Here the factor $3$ comes from $\log D_n=n+o(n)$; dropping it would invalidate the arithmetic comparison. These are properties of this particular unreduced integral family, not optimal irrationality-measure bounds \eref{1283}{1}, \eref{1283}{2}.

For an algebraic irrational $\alpha$ of degree $d$, elementary polynomial evaluation at $B/A$ gives
\begin{equation}\label{eq184:liouville}
 |A\alpha-B|\ge c_\alpha\max(|A|,|B|)^{1-d}
 \quad(A\ne0),
\end{equation}
The inequality follows by clearing $A^d$ from a minimal polynomial. More explicitly, a nonzero integer $A^dP(B/A)$ has modulus at least one, and $|P(B/A)|\le C|B/A-\alpha|$ when $B/A$ is close to $\alpha$; ratios bounded away from $\alpha$ give a uniform lower bound directly (for unbounded ratios, $|A\alpha-B|$ is comparable to $|B|$).

For the displayed Ap\'ery family the small-value-to-height exponent is
\[
 \tau=\frac{\log\rho-3}{\log\rho+3}<1.
\]
It therefore does not contradict \eqref{eq184:liouville} even for $d=2$. The proposed tensor repair $Q_{n,r}(T)=(A_nT-B_n)^r$ has
\[
 \log H(Q_{n,r})=r(3+\log\rho)n+o(n),\qquad
 \log|Q_{n,r}(\zeta(3))|=-r(\log\rho-3)n+o(n)
\]
for each fixed $r$. The exponent ratio is unchanged. More decisively, evaluating a single such power is logically just evaluating the original linear form: it creates no independent relation. Products of finitely many comparable approximants have the same defect unless an additional coefficient-cancellation or determinant mechanism is proved.

\paragraph{Stress test.}
The enclosure is conditional only on the established identification of the recurrence limit with $\zeta(3)$; it does not reprove that identification. Integrality was retained before comparing height and error. We do not exclude improvements arising from unexpected common divisors of $A_n,B_n$, optimized linear forms, or genuinely independent Hermite--Pad\'e systems. We exclude only the claimed gain obtained by formally taking powers of the displayed family. The Liouville comparison is a diagnostic, not an assertion that it is the only possible transcendence method. A large finite collection of nonzero polynomial evaluations cannot replace the unbounded degree and height quantifiers.

\paragraph{Outcome.}
\textbf{Outcome: Exploratory approach with unresolved gap.}
An exact rational enclosure and a quantitative obstruction to the simple tensor-power proposal were obtained. The recurrence identities are classical in origin and no literature priority is claimed. A complete resolution would require genuinely new higher-degree arithmetic information, for example independent small polynomial forms with a proved nonvanishing determinant and a height/error balance strong enough for each algebraic degree. None is supplied here.

% END RESEARCH ADDITION ATTEMPT-184
\subsection{Sources}
\begin{itemize}\raggedright
\item[\textnormal{[S1]}]\hypertarget{source:1283:1}{} W. Zudilin, Arithmetic of linear forms involving odd zeta values, J. Theor. Nombres Bordeaux 16 (2004).
\url{https://arxiv.org/abs/math/0206176}.
{\small\textit{Catalogue formulation source.}}
% BEGIN RESEARCH ADDITION REFERENCES-184
\item[\textnormal{[E1]}]\hypertarget{extra:1283:1}{} Henri Cohen and Wadim Zudilin, \emph{Variations on a theme of Ap\'ery}, arXiv:2501.10090v1 (17 January 2025), Section 1: Ap\'ery recurrence, rational approximants and their analytic behavior.
\url{https://arxiv.org/html/2501.10090v1}.
{\small\textit{Research reference; consulted 27 September 2026.}}
\item[\textnormal{[E2]}]\hypertarget{extra:1283:2}{} Wadim Zudilin, \emph{An elementary proof of Ap\'ery\textquotesingle s theorem}, arXiv:math/0202159; comparison of linear-form decay with denominator growth.
\url{https://arxiv.org/abs/math/0202159}.
{\small\textit{Research reference; consulted 27 September 2026.}}
% END RESEARCH ADDITION REFERENCES-184
\end{itemize}

\end{document}
